\chapter{Operaciones intrínsecas correlativas del mismo continuo}
\label{chap:pdfv2-operaciones-intrinsecas-correlativas}

\begin{thesisbox}
Los símbolos \(\mathrm{sp},\mathrm{sol},\mathrm{gau},\mathrm{coh}\) y
\(\mathrm{det}\) que aparecen en este capítulo marcan cinco operaciones
correlativas sobre un único estado APP--TRIT--TPK\@. Cada operación recibe
la misma extensión superviviente, la
misma orientación, el mismo carry, la misma memoria, la misma frontera y la
misma ruta. Los apartados siguientes son etapas matemáticas sucesivas de una
sola construcción y no exposiciones autónomas.
\end{thesisbox}

\section{Arista generadora y estado genealógico común}

No repetimos aquí el estado, el árbol ni la medida construidos en el capítulo
anterior. Tomamos
\(\widehat x=\widehat x_k(\gamma)\in
\widehat{\mathcal X}^{\rm tot}_k\) de
\eqref{eq:continuo-comun-estado} y todas sus emisiones admisibles
\(\alpha\in\operatorname{Out}(\widehat x)\). Si
\(\widehat y=\alpha_*\widehat x\), escribimos
\begin{equation}
\label{eq:pdfv2-cor-hijos-supervivientes}
 \operatorname{Ch}_k(\widehat x):=
 \{(\alpha,\alpha_*\widehat x):
      \alpha\in\operatorname{Out}(\widehat x)\},
 \qquad
 d_k(\widehat x):=\#\operatorname{Ch}_k(\widehat x).
\end{equation}
La finitud y la no vacuidad de este conjunto son ya
\eqref{eq:continuo-comun-out-finito}; la supervivencia a todos los horizontes,
la multisección completa y su medida son
\cref{thm:continuo-comun-total-superviviente,thm:continuo-comun-multiseccion-no-vacia,thm:continuo-comun-medida-canonica}. No se vuelve a suponer ni a probar una
segunda no vacuidad.

Cada par \((\alpha,\widehat y)\) de
\eqref{eq:pdfv2-cor-hijos-supervivientes} conserva el registro único
\begin{equation}
\label{eq:pdfv2-cor-registro-extension}
 \begin{aligned}
 \mathfrak e_\alpha:=\bigl(&
 \alpha,\operatorname{Ext}_{\gamma_\alpha};
 \operatorname{Sh}_{s(\alpha)},\operatorname{Sh}_{t(\alpha)};\\
 &
 \operatorname{res}^{\Sigma,\Pi}(\alpha),
 \operatorname{quo}^{\Sigma,\Pi}(\alpha);\\
 &\operatorname{or}(\alpha),\kappa(\alpha),\mathsf C_\alpha;
 \mathsf Q(s(\alpha)),\mathsf Q(t(\alpha)),g(s(\alpha)),g(t(\alpha)),
 q^{(9)}(s(\alpha)),q^{(9)}(t(\alpha));\\
 &
 m(\gamma\alpha;m_0),\partial\alpha,\gamma\alpha,
 \operatorname{Led}(\gamma\alpha);\\
 &\{\operatorname{Term}_{k+1,N}(\widehat y)\}_{N\geq k+1},
 \operatorname{Msect}(\widehat y),\mu_{\widehat y}\bigr).
 \end{aligned}
\end{equation}
Aquí \(\operatorname{Ext}_{\gamma_\alpha}\) nombra la relación completa que
la emisión admisible ya satisface; no introduce otro filtro. Las operaciones
posteriores son aplicaciones coordinadas de \(\mathfrak e_\alpha\) y no
reconstruyen el registro desde sus salidas.

\section{Acciones intrínsecas correlativas sobre el generador común}

Las marcas internas se fijan sin cambiar de fuente:
\(\mathrm{sp}\) designa transporte y graduación de las dos hojas;
\(\mathrm{sol}\), incremento firmado, balance y memoria;
\(\mathrm{gau}\), incidencia y curvatura de carry;
\(\mathrm{coh}\), puertos, diferencia rectangular y complejo celular; y
\(\mathrm{det}\), complejo basado, relleno correlativo y línea determinante. Las fórmulas
siguientes muestran sus acoplamientos sobre el mismo
\(\mathfrak e_\alpha\); sus acciones comparten literalmente fuente, destino,
ruta y ledger.

La doble evaluación APP proporciona, sobre el mismo soporte, las etiquetas
\(\Sigma\) y \(\Pi\). Para una extensión \(\alpha\), sus multiplicidades
se calculan del ledger completo, sin un contador exterior:
\begin{equation}
\label{eq:pdfv2-cor-multiplicidades-hojas}
 N_\diamond(\alpha):=
 \sum_{\epsilon\ {\rm en}\ \gamma\alpha}
 \mathbf1_{\{\operatorname{res}^{\diamond}(\epsilon)\neq0\}},
 \qquad \diamond\in\{\Sigma,\Pi\}.
\end{equation}
Sea además el coeficiente orientado total
\[
 \Delta_{\gamma_\alpha}:=
 \sum_{\epsilon\ {\rm en}\ \gamma\alpha}\operatorname{or}(\epsilon)
 \in\mathbb Z.
\]
La acción firmada es
\begin{equation}
\label{eq:pdfv2-cor-b-alpha}
 b_\alpha:=\Delta_{\gamma_\alpha}
 \bigl(N_\Sigma(\alpha)-N_\Pi(\alpha)\bigr),
 \qquad
 b_\alpha^\pm:=\frac{|b_\alpha|\pm b_\alpha}{2}.
\end{equation}
Por construcción,
\begin{equation}
\label{eq:pdfv2-cor-b-identidades}
 b_\alpha=b_\alpha^+-b_\alpha^-,\qquad
 |b_\alpha|=b_\alpha^++b_\alpha^-,\qquad
 b_\alpha^+b_\alpha^-=0.
\end{equation}

El mismo conjunto de extensiones produce simultáneamente los puertos
\begin{equation}
\label{eq:pdfv2-cor-puertos-comunes}
 I_k(\widehat x_k):=
 \operatorname{Ch}_k(\widehat x_k)\times\{\Sigma\},
 \qquad
 J_k(\widehat x_k):=
 \operatorname{Ch}_k(\widehat x_k)\times\{\Pi\}.
\end{equation}
No se supone que una hoja visible contenga más prolongaciones que la otra:
cada extensión superviviente se evalúa en las dos hojas del mismo soporte.
Por \eqref{eq:continuo-comun-out-finito}, ambos conjuntos de puertos son no
vacíos.

Para \(A_N=\mathbb Z/3^N\mathbb Z\), la acción celular sobre generadores es
\begin{equation}
\label{eq:pdfv2-cor-kappa-celda}
 \begin{aligned}
 \kappa_{k,N}:A_N^{I_k}\oplus A_N^{J_k}
 &\longrightarrow A_N^{I_k\times J_k},\\
 \kappa_{k,N}(u,v)_{(\alpha,\Sigma),(\beta,\Pi)}
 &:=u_{(\alpha,\Sigma)}-v_{(\beta,\Pi)}.
 \end{aligned}
\end{equation}
Elegidos los primeros puertos según el orden TPK, la diferencia rectangular
\begin{equation}
\label{eq:pdfv2-cor-delta-celda}
 \delta(w)_{ij}:=w_{ij}-w_{ij_0}-w_{i_0j}+w_{i_0j_0}
\end{equation}
descarga la sucesión escindida
\begin{equation}
\label{eq:pdfv2-cor-sucesion-celda}
 0\longrightarrow A_N\xrightarrow{\Delta}
 A_N^{I_k}\oplus A_N^{J_k}
 \xrightarrow{\kappa_{k,N}}A_N^{I_k\times J_k}
 \xrightarrow{\delta}A_N^{(|I_k|-1)(|J_k|-1)}
 \longrightarrow0.
\end{equation}
Ésta es una identidad de la misma fibra APP: no añade una celda ajena al
generador \(\mathfrak e_\alpha\).

La acción celular local usa las dos coordenadas enteras acumuladas por ese
mismo ledger. Para
\(\widehat x=\widehat x_k(\gamma)\) y
\(\widehat y=\alpha_*\widehat x=\widehat x_{k+1}(\gamma\alpha)\), ponemos
\begin{equation}
\label{eq:pdfv2-cor-c-v}
 \begin{aligned}
 c_{\widehat x}&:=\kappa(\gamma), &
 c_\alpha:=c_{\widehat y}&:=\kappa(\gamma\alpha),\\
 v_{\widehat x}&:=\sum_{\epsilon\ {\rm en}\ \gamma}
                  \operatorname{or}(\epsilon), &
 v_{\widehat y}&:=\sum_{\epsilon\ {\rm en}\ \gamma\alpha}
                  \operatorname{or}(\epsilon),\\
 v_\alpha&:=v_{\widehat y}-v_{\widehat x}
           =\operatorname{or}(\alpha).&&
 \end{aligned}
\end{equation}
Sobre los generadores \(f,r,e,b,g\) se define
\begin{equation}
\label{eq:pdfv2-cor-complejo-local}
 \partial f=3c_\alpha e+b,\qquad
 \partial r=0,\qquad
 \partial e=g,\qquad
 \partial b=-3c_\alpha g,\qquad
 \partial g=0.
\end{equation}
La cancelación
\begin{equation}
\label{eq:pdfv2-cor-dos-cero}
 \partial^2f=3c_\alpha g-3c_\alpha g=0
\end{equation}
es literal y usa el mismo carry con los dos signos impuestos por la
orientación.

La semilla de incidencia se obtiene del módulo libre de las dos hojas,
\begin{equation}
\label{eq:pdfv2-cor-v-tpk}
 V_{\rm TPK}:=\mathbf F_3e_\Sigma\oplus\mathbf F_3e_\Pi,
 \qquad
 \mathscr I:=\mathbf P^1(V_{\rm TPK}),
 \qquad |\mathscr I|=4,
\end{equation}
y de sus pares no ordenados
\begin{equation}
\label{eq:pdfv2-cor-incidencias}
 \mathscr E:=\binom{\mathscr I}{2},
 \qquad |\mathscr E|=6,
 \qquad
 \mathscr B:=\mathscr I\times\{+,-\},
 \qquad |\mathscr B|=8.
\end{equation}
La aplicación de incidencia queda determinada en la base por
\begin{equation}
\label{eq:pdfv2-cor-b-incidencia-generadores}
 B[e_L]=\sum_{L'\neq L}e_{\{L,L'\}},
 \qquad L\in\mathscr I,
\end{equation}
y, por linealidad,
\begin{equation}
\label{eq:pdfv2-cor-b-incidencia-coordenadas}
 (Ba)_{\{L,L'\}}=a_L+a_{L'}.
\end{equation}
Si \(s(q)=\sum_{E\in\mathscr E}q_E\), cada \(L\) aparece tres veces y
\begin{equation}
\label{eq:pdfv2-cor-sb-cero}
 s(Ba)=0.
\end{equation}
Por tanto la función
\begin{equation}
\label{eq:pdfv2-cor-wc}
 W_c(q):=\mathbf1_{\{s(q)=0\}}-\frac13
\end{equation}
satisface \(W_c(q+Ba)=W_c(q)\). Las identidades
\eqref{eq:pdfv2-cor-v-tpk}--\eqref{eq:pdfv2-cor-wc} son internas y exactas.
El lift material de \(q\) desde las celdas y los carries comunes se trata en
la capa de ensamblaje; no se confunde esta identidad de incidencia con la
construcción de ese lift.

\section{Transporte unitario, polarización y balance sobre la fibra común}

La no vacuidad de \eqref{eq:pdfv2-cor-hijos-supervivientes} permite fijar,
sin seleccionar una prolongación,
\begin{equation}
\label{eq:pdfv2-cor-peso-conteo}
 p_k(\alpha,\widehat y\mid\widehat x)
 :=\frac1{d_k(\widehat x)}.
\end{equation}
Sea el espacio de las dos hojas del mismo estado
\begin{equation}
\label{eq:pdfv2-cor-hilbert-hojas}
 \mathscr H_k:=
 \bigoplus_{\widehat x\in\widehat{\mathcal X}^{\rm tot}_k}
 \bigl(\mathsf H_{\widehat x}^{\Sigma}
       \oplus\mathsf H_{\widehat x}^{\Pi}\bigr),
\end{equation}
y sea \(\iota_{\widehat x}\) la inclusión de la fibra indicada. El transporte
TPK de \eqref{eq:continuo-comun-transporte-hojas} da la acción exacta
\begin{equation}
\label{eq:pdfv2-cor-u-comun}
 U_k\iota_{\widehat x}v:=
 \sum_{(\alpha,\widehat y)\in\operatorname{Ch}_k(\widehat x)}
 p_k(\alpha,\widehat y\mid\widehat x)^{1/2}
 \iota_{\widehat y}U_\alpha v.
\end{equation}
Los hijos de dos padres distintos son disjuntos porque el truncamiento es una
función. Además, dos testigos \(\alpha\) distintos no colapsan en el mismo
sumando directo: la ruta completa \(\gamma\), incluido su último paso, es una
coordenada de \(\widehat y\). En consecuencia, la identidad de Gram es
\begin{equation}
\label{eq:pdfv2-cor-u-gram}
 \left.U_k^*U_k\right|_{\mathsf H_{\widehat x}^{\Sigma}
       \oplus\mathsf H_{\widehat x}^{\Pi}}
 =
 \sum_{(\alpha,\widehat y)\in\operatorname{Ch}_k(\widehat x)}
 p_k(\alpha,\widehat y\mid\widehat x)\,U_\alpha^*U_\alpha.
\end{equation}
Por tanto \(U_k^*U_k=I\) exactamente cuando los transportes de hoja
\(U_\alpha\) publicados por TPK son isometrías. La normalización de los pesos,
por sí sola, no puede sustituir esa comprobación.
Para \(\ell>k\) se conserva la historia completa mediante
\begin{equation}
\label{eq:pdfv2-cor-u-iterado}
 U_{\ell,k}:=U_{\ell-1}\cdots U_k,
 \qquad U_{k,k}:=I.
\end{equation}

La graduación actúa sobre las dos coordenadas que ya están presentes en cada
fibra de \eqref{eq:pdfv2-cor-hilbert-hojas}:
\begin{equation}
\label{eq:pdfv2-cor-sigma}
 \sigma_k\iota_{\widehat x}(v_\Sigma,v_\Pi)
 :=\iota_{\widehat x}(v_\Sigma,-v_\Pi),
\qquad
 P_k:=\frac{I+\sigma_k}{2},\quad
 Q_k:=\frac{I-\sigma_k}{2}.
\end{equation}
El umbral TRIT pertenece a la orientación y al carry de
\(\widehat x\); no borra ninguna de las dos hojas APP y no requiere un tercer
caso en \eqref{eq:pdfv2-cor-sigma}.

La parte que conserva hoja y la parte que la cambia son dos componentes del
mismo \(U_k\):
\begin{align}
 A_k&:=P_{k+1}U_kP_k+Q_{k+1}U_kQ_k,
 \label{eq:pdfv2-cor-a-comun}\\
 \eta_k&:=Q_{k+1}U_kP_k+P_{k+1}U_kQ_k.
 \label{eq:pdfv2-cor-eta-comun}
\end{align}
Por ortogonalidad de \(P_{k+1}\) y \(Q_{k+1}\),
\begin{equation}
\label{eq:pdfv2-cor-balance-hojas}
 U_k=A_k+\eta_k,\qquad
 A_k^*A_k+\eta_k^*\eta_k=U_k^*U_k.
\end{equation}
El último miembro es \(I\) bajo, y sólo bajo, la condición isométrica
identificada después de \eqref{eq:pdfv2-cor-u-gram}.
Esta identidad y el incremento \eqref{eq:pdfv2-cor-b-alpha} no son dos
transportes: \(A_k,\eta_k\) actúan sobre la base de historias y
\(b_\alpha\) lee la arista de esa misma base.

\subsection{Polarización firmada, telescopía no autónoma y retorno sin reinicialización}
No se introduce aquí una rama espectral independiente. Se restringen al
mismo espacio \(\mathscr H_k\) y al mismo transporte \(U_k\) ya construidos
las notaciones
\[
 \mathcal H_k^{\mathrm{sp}}:=\mathscr H_k,
 \qquad
 \Lambda(\gamma_{\widehat x})
 :=\bigl(\operatorname{or},\operatorname{Carr},\operatorname{Mem},
          \partial,\gamma\bigr)(\widehat x),
\]
y \(\Sigma_{\rm reg}\Lambda(\gamma_{\widehat x})\) designa su familia de
truncamientos compatibles. De este modo, el lector firmado, el defecto y el
terminal siguientes son acciones del constructor correlativo común, no una
selección retrospectiva de historias.
\begingroup
  \let\section\subsection
  \input{incorporaciones/v2_global_20260823/05a_lector_telescopia_correlativa}
  \input{incorporaciones/v2_global_20260823/05a_retorno_memoria_correlativa}
\endgroup

La esperanza condicionada asociada al mismo peso es
\begin{equation}
\label{eq:pdfv2-cor-esperanza-comun}
 (E_kF)(\widehat x):=
 \sum_{(\alpha,\widehat y)\in\operatorname{Ch}_k(\widehat x)}
 p_k(\alpha,\widehat y\mid\widehat x)F(\alpha,\widehat y).
\end{equation}
Para el incremento fino \(b_f\) se ponen
\begin{equation}
\label{eq:pdfv2-cor-kappa-jensen}
 b_c:=E_kb_f,\qquad
 \kappa_k^{\rm can}:=\frac{E_k|b_f|-|b_c|}{2}\geq0.
\end{equation}
Entonces
\begin{equation}
\label{eq:pdfv2-cor-jensen-hojas}
 E_kb_f^+=b_c^++\kappa_k^{\rm can},\qquad
 E_kb_f^-=b_c^-+\kappa_k^{\rm can}.
\end{equation}
El superíndice \({\rm can}\) distingue este carry de cancelación del carry
entero \(\kappa_\alpha\); ambos proceden del mismo transporte, pero tienen
tipos distintos.

Las memorias firmadas no requieren una condición inicial exterior: para
\(\gamma=(\epsilon_1,\ldots,\epsilon_k)\) se definen por suma del propio
ledger,
\begin{equation}
\label{eq:pdfv2-cor-memorias}
 \begin{aligned}
 M_0^\pm&:=0,\\
 M_k^\pm(\gamma)&:=\sum_{j=1}^k b_{\epsilon_j}^\pm,\\
 M_{k+1}^\pm(\gamma\alpha)&:=M_k^\pm(\gamma)+b_\alpha^\pm,\\
 D_k^0&:=M_k^+-M_k^-, &
 H_k^0&:=M_k^++M_k^-.
 \end{aligned}
\end{equation}
Así, el balance firmado, la memoria total y la polarización
\(\langle v,\sigma_kv\rangle\) son tres lecturas coordinadas del mismo
prefijo, no tres estados reconstruidos por separado.

\subsection{Refinamiento nonádico, pérdidas ortogonales y memoria terminal}
El índice \(j\) del desarrollo siguiente es el nivel del mismo árbol de
prolongaciones y su registro de trece coordenadas es la restricción del
registro común de \eqref{eq:pdfv2-cor-registro-extension}. Por tanto se fija
\(D_{{\rm sol},j}:=D_j^{(13)}\) y no se postula un dominio adicional.
\begingroup
  \let\section\subsection
  \input{incorporaciones/v2_global_20260823/05b_balance_memoria_correlativa}
\endgroup

\section{Orientación, acarreo y composición simultánea de trayectorias}

La inversión de la ruta actúa una sola vez sobre
\(\mathfrak e_\alpha\). Sea
\(\mathsf J(v_\Sigma,v_\Pi):=(v_\Pi,v_\Sigma)\) el intercambio de hojas,
y sea \(P_{\rm tr}\) la transposición de matrices de puertos. Sus cinco
expresiones locales son
\begin{equation}
\label{eq:pdfv2-cor-orientacion-cruzada}
 \begin{gathered}
 \mathsf J\sigma_k=-\sigma_k\mathsf J,
 \qquad b_{\bar\alpha}=-b_\alpha,
 \qquad (b_{\bar\alpha}^+,b_{\bar\alpha}^-)
       =(b_\alpha^-,b_\alpha^+),\\
 \vartheta_L(K,\epsilon):=
 \begin{cases}
  (L,-\epsilon),&K=L,\\
  (K,\epsilon),&K\neq L,
 \end{cases}
 \qquad
 \kappa_{J,I}(v,u)=-P_{\rm tr}\kappa_{I,J}(u,v),\\
 \kappa(\bar\alpha)
 =-\mathsf C_\alpha^{-1}\kappa(\alpha),
 \qquad v_{\bar\alpha}=-v_\alpha.
 \end{gathered}
\end{equation}
Aquí \(\mathsf J\) intercambia las dos hojas, \(P_{\rm tr}\) transpone la matriz de
incidencia y \(\vartheta_L\) intercambia exactamente las dos caras de la
dirección \(L\), fijando las otras seis. En la base libre de
\(\mathscr B=\mathscr I\times\{+,-\}\), los ocho vectores son distintos,
\(\vartheta_L^2=I\) y
\(\vartheta_L\neq\vartheta_{L'}\) para \(L\neq L'\). La misma orientación
explica todos los signos de
\eqref{eq:pdfv2-cor-orientacion-cruzada}.

Para extensiones componibles \(\alpha:x\to y\) y \(\beta:y\to z\), el
carry nativo satisface
\begin{equation}
\label{eq:pdfv2-cor-carry-nativo}
 \kappa(\beta\alpha)=
 \kappa(\beta)+\mathsf C(\beta)\kappa(\alpha).
\end{equation}
La elevación matricial no se deja como un nombre. En la base ordenada
\((f,r,e,b,g)\), sea
\(\widetilde L_N(\alpha)\in\operatorname{Mat}_5(\mathbb Z)\) el
representante entero centrado de la matriz de
\(\Psi_\alpha\) de \eqref{eq:pdfv2-cor-psi-generadores} módulo \(3^N\), y
defínase
\begin{equation}
\label{eq:pdfv2-cor-carry-matricial-def}
 c_N(\beta,\alpha):=
 \frac{\widetilde L_N(\beta)\widetilde L_N(\alpha)
       -\widetilde L_N(\beta\alpha)}{3^N}
 \in\operatorname{Mat}_5(\mathbb Z).
\end{equation}
La divisibilidad procede de la igualdad de las dos composiciones módulo
\(3^N\). Escribiendo \(L_N=[\widetilde L_N]_{3^N}\), los lifts enteros
satisfacen
\begin{equation}
\label{eq:pdfv2-cor-carry-matricial}
 \widetilde L_N(\beta)\widetilde L_N(\alpha)
 =\widetilde L_N(\beta\alpha)+3^Nc_N(\beta,\alpha)
\end{equation}
y la asociatividad produce el cociclo
\begin{equation}
\label{eq:pdfv2-cor-cociclo-matricial}
 c_N(\delta,\beta\alpha)+\widetilde L_N(\delta)c_N(\beta,\alpha)
 =c_N(\delta\beta,\alpha)
  +c_N(\delta,\beta)\widetilde L_N(\alpha).
\end{equation}

El mismo incremento de carry actúa sobre el complejo local mediante
\begin{equation}
\label{eq:pdfv2-cor-psi-generadores}
 \begin{gathered}
 \Psi_\alpha(r_x)=r_y,\quad
 \Psi_\alpha(e_x)=e_y,\quad
 \Psi_\alpha(g_x)=g_y,\quad
 \Psi_\alpha(f_x)=f_y,\\
 \Psi_\alpha(b_x)=b_y+3(c_y-c_x)e_y.
 \end{gathered}
\end{equation}
El cambio de frontera orientada queda en
\begin{equation}
\label{eq:pdfv2-cor-k-alpha}
 K_\alpha:=(v_y-v_x)f_y,
 \qquad
 K_{\beta\alpha}=K_\beta+\Psi_\beta K_\alpha.
\end{equation}
Las ecuaciones \eqref{eq:pdfv2-cor-carry-nativo},
\eqref{eq:pdfv2-cor-cociclo-matricial} y
\eqref{eq:pdfv2-cor-k-alpha} son tres grados del mismo ledger de composición.
Ninguna autoriza a borrar las otras dos.

\subsection{Nervio de prolongaciones y diagonal de Alexander--Whitney}
Para una historia del registro común se escribe
\(A_n:=\mathbf Z/3^n\mathbf Z\) y \(\mathsf T_m(x)\) para la categoría
finita de \emph{todos} sus prefijos TPK de profundidad \(m\), con los
morfismos de extensión ya construidos. La diagonal y la counidad que siguen
se aplican a ese mismo nervio; no definen una quinta rama autónoma.
\begingroup
  \let\section\subsection
  \input{incorporaciones/v2_global_20260823/05e_aw_nervio_correlativo}
\endgroup

\section{Incidencia celular y relleno correlativo de una misma trayectoria}

La sucesión \eqref{eq:pdfv2-cor-sucesion-celda} asigna a cada prefijo el
complejo de dos términos
\begin{equation}
\label{eq:pdfv2-cor-cono-celda}
 C_{k,N}(\widehat x):=
 \left[A_N^{I_k}\oplus A_N^{J_k}
 \xrightarrow{\kappa_{k,N}}A_N^{I_k\times J_k}\right].
\end{equation}
Un refinamiento actúa en los puertos por truncamiento de la extensión y en
los coeficientes por \(A_{N'}\twoheadrightarrow A_N\). Cuando
\(\alpha_u:I_{k'}\to I_k\) y \(\beta_u:J_{k'}\to J_k\) son esos mapas, su
acción en generadores es
\begin{equation}
\label{eq:pdfv2-cor-refinamiento-celda}
 (p,q)\longmapsto(p\circ\alpha_u,q\circ\beta_u),
 \qquad
 w\longmapsto w\circ(\alpha_u\times\beta_u),
\end{equation}
y verifica
\begin{equation}
\label{eq:pdfv2-cor-refinamiento-cadena}
 \kappa_{k',N}\Gamma(u)_1=\Gamma(u)_0\kappa_{k,N}.
\end{equation}

La incidencia no se fabrica etiquetando cuatro copias de una emisión neutra.
Se genera por la completación celular libre del módulo de las dos hojas que
ya porta el mismo estado. Para
\[
 V_{\widehat x}:=\mathbf F_3e_{\Sigma,\widehat x}
 \oplus\mathbf F_3e_{\Pi,\widehat x},
\]
las cuatro rectas internas son
\begin{equation}
\label{eq:pdfv2-cor-cuatro-rectas-tpk}
 \mathscr I_{\widehat x}:=
 \{L_0,L_\infty,L_+,L_-\}
 =\mathbf P^1(V_{\widehat x}),\qquad
 \begin{aligned}
 L_0&=\langle e_\Sigma\rangle,&
 L_\infty&=\langle e_\Pi\rangle,\\
 L_+&=\langle e_\Sigma+e_\Pi\rangle,&
 L_-&=\langle e_\Sigma-e_\Pi\rangle.
 \end{aligned}
\end{equation}
Si \(v_L\) es, respectivamente,
\((1,0)^t,(0,1)^t,(1,1)^t,(1,-1)^t\), la acción de la dirección \(L\)
sobre los generadores de hoja es el proyector
\begin{equation}
\label{eq:pdfv2-cor-proyectores-cuatro-direcciones}
 P_L:=v_L(v_L^tv_L)^{-1}v_L^t,
\end{equation}
es decir, sobre \(\mathbf F_3\),
\[
 P_0=\begin{pmatrix}1&0\\0&0\end{pmatrix},\quad
 P_\infty=\begin{pmatrix}0&0\\0&1\end{pmatrix},\quad
 P_+=\begin{pmatrix}2&2\\2&2\end{pmatrix},\quad
 P_-=\begin{pmatrix}2&1\\1&2\end{pmatrix}.
\]
El cálculo directo da
\begin{equation}
\label{eq:pdfv2-cor-proyectores-no-colapso}
 P_L^2=P_L,\qquad \operatorname{rk}P_L=1,\qquad
 \operatorname{Im}P_L=L,\qquad
 P_L=P_{L'}\Longrightarrow L=L'.
\end{equation}
Por tanto las cuatro direcciones son salidas distintas del módulo
\(V_{\widehat x}\); no son cuatro nombres puestos sobre la misma acción.

\begin{definition}[Completación celular libre de las hojas TPK]
\label{def:pdfv2-cor-completacion-celular-tpk}
Sea \(\operatorname{Cell}_{\rm TPK}(\widehat x)\) el complejo libre
\[
 C_2(\widehat x)\xrightarrow{\partial_2}
 C_1(\widehat x)\xrightarrow{\partial_1}C_0(\widehat x)
\]
con
\begin{align*}
 C_0(\widehat x)
 &:=
 \mathbf F_3[v_\star]\oplus
 \bigoplus_{L\in\mathscr I}\mathbf F_3[v_L]\oplus
 \bigoplus_{(L,\epsilon)\in\mathscr B}
       \mathbf F_3[v_{L,\epsilon}]\oplus
 \bigoplus_{E\in\mathscr E}\mathbf F_3[v_E],\\
 C_1(\widehat x)
 &:=
 \bigoplus_{L\in\mathscr I}\mathbf F_3[u_L]\oplus
 \bigoplus_{(L,\epsilon)\in\mathscr B}
       \mathbf F_3[u_{L,\epsilon}]\oplus
 \bigoplus_{E=\{L,L'\}\in\mathscr E}
 \bigl(\mathbf F_3[u_{L'\mid L}]
       \oplus\mathbf F_3[u_{L\mid L'}]\bigr),\\
 C_2(\widehat x)
 &:=
 \bigoplus_{E\in\mathscr E}\mathbf F_3[\omega_E].
\end{align*}
Su acción en generadores es
\begin{align}
 \partial u_L&=v_L-v_\star,
 &
 \partial u_{L'\mid L}&=v_E-v_L,
 \label{eq:pdfv2-cor-cell-frontera-aristas}\\
 \partial u_{L,\epsilon}&=v_{L,\epsilon}-v_L,
 &
 \partial\omega_{\{L,L'\}}
 &=u_L+u_{L'\mid L}-u_{L'}-u_{L\mid L'}.
 \label{eq:pdfv2-cor-cell-frontera-celda}
\end{align}
\end{definition}

Los ocho generadores \(v_{L,\epsilon}\) realizan materialmente la pantalla
\(\mathscr B\). La involución anterior se extiende al complejo por
\begin{equation}
\label{eq:pdfv2-cor-pantalla-ocho-generadores}
 \vartheta_Lv_{K,\epsilon}:=v_{\vartheta_L(K,\epsilon)},\qquad
 \vartheta_Lu_{K,\epsilon}:=u_{\vartheta_L(K,\epsilon)},
\end{equation}
y fija los demás generadores. Por ser libre la base,
\begin{equation}
\label{eq:pdfv2-cor-pantalla-ocho-no-colapso}
 |\{v_{L,\epsilon}:(L,\epsilon)\in\mathscr B\}|=8,\qquad
 \vartheta_L^2=I,\qquad
 \vartheta_L\vartheta_{L'}=\vartheta_{L'}\vartheta_L.
\end{equation}
Además \(\partial\vartheta_L=\vartheta_L\partial\), pues ambos recorridos
envían \(u_{L,\epsilon}\) a \(v_{L,-\epsilon}-v_L\) y fijan las fronteras
restantes.

Las dos palabras
\[
 p_{L,L'}:=u_{L'\mid L}u_L,\qquad
 p_{L',L}:=u_{L\mid L'}u_{L'}
\]
son caminos diferentes del complejo libre. Su diferencia es la frontera de
\(\omega_{\{L,L'\}}\), y
\begin{equation}
\label{eq:pdfv2-cor-cell-dos-cero}
 \partial^2\omega_E
 =(v_L-v_\star)+(v_E-v_L)
 -(v_{L'}-v_\star)-(v_E-v_{L'})=0.
\end{equation}
La representación de TPK en las hojas se fija por
\[
 \varrho(u_L)=P_L,\qquad
 \varrho(u_{L'\mid L})=P_{L'},\qquad
 \varrho(p_{L,L'})=P_{L'}P_L.
\]
Aunque dos matrices publicadas coincidieran en una especialización, los dos
caminos no colapsan: siguen siendo generadores libres distintos de
\(C_1(\widehat x)\), unidos por una 2--celda explícita.

La coordenada completa de incidencia es el grupo de 2--cocadenas
\begin{equation}
\label{eq:pdfv2-cor-q-incidencia-total}
 Q_{\rm inc}(\widehat x):=
 C^2(\operatorname{Cell}_{\rm TPK}(\widehat x);\mathbf F_3)
 =\bigoplus_{E\in\mathscr E}\mathbf F_3\omega_E^*
 \simeq\mathbf F_3^6.
\end{equation}
La operación interna devuelve todo \(Q_{\rm inc}(\widehat x)\), no un
elemento seleccionado. Cada \(q\) tiene expansión única
\(q=\sum_Eq_E\omega_E^*\). Para
\(a\in\mathbf F_3^{\mathscr I}\), su cambio de potencial es
\begin{equation}
\label{eq:pdfv2-cor-collar-variacion}
 q\longmapsto q+Ba,\qquad (Ba)_{\{L,L'\}}=a_L+a_{L'}.
\end{equation}
Como cada recta pertenece a tres pares,
\[
 s(Ba)=3\sum_{L\in\mathscr I}a_L=0,
\]
y por ello
\begin{equation}
\label{eq:pdfv2-cor-wc-invariante-celular}
 W_c(q+Ba)=W_c(q)\qquad(q\in Q_{\rm inc}(\widehat x)).
\end{equation}
La unidad de la completación sólo ancla el complejo en el estado:
\begin{equation}
\label{eq:pdfv2-cor-cell-unidad}
 \eta^{\rm cell}_{\widehat x}:V_{\widehat x}\longrightarrow
 V_{\widehat x}\otimes C_0(\operatorname{Cell}_{\rm TPK}(\widehat x)),\qquad
 e_{\diamond,\widehat x}\longmapsto
 e_{\diamond,\widehat x}\otimes v_\star,
\end{equation}
y marca el origen \(q=0\); no afirma que la órbita de ese origen sea todo
\(Q_{\rm inc}\).

La construcción es funtorial en el truncamiento común. Si
\(\tau_{\ell,k}\widehat x_\ell=\widehat x_k\), definimos
\begin{align*}
 \tau_V&:V_{\widehat x_\ell}\longrightarrow V_{\widehat x_k},&
 \tau_Ve_{\diamond,\widehat x_\ell}&=e_{\diamond,\widehat x_k},\\
 \tau_a&:\mathbf F_3^{\mathscr I_{\widehat x_\ell}}
          \longrightarrow\mathbf F_3^{\mathscr I_{\widehat x_k}},&
 (\tau_aa)_{\tau_{\mathscr I}L}&=a_L,\\
 \tau_Q&:Q_{\rm inc}(\widehat x_\ell)
          \longrightarrow Q_{\rm inc}(\widehat x_k),&
 (\tau_Qq)_{\tau_{\mathscr E}E}&=q_E,
\end{align*}
donde \(\tau_{\mathscr I}[v]=[\tau_Vv]\) y
\(\tau_{\mathscr E}\{L,L'\}=\{\tau_{\mathscr I}L,
\tau_{\mathscr I}L'\}\). El mapa de cadenas
\(\operatorname{Cell}(\tau):C_\bullet(\widehat x_\ell)
\to C_\bullet(\widehat x_k)\) actúa en generadores por
\begin{equation}
\label{eq:pdfv2-cor-cell-truncamiento-generadores}
 v_\bullet(\widehat x_\ell)\mapsto v_\bullet(\widehat x_k),\quad
 u_\bullet(\widehat x_\ell)\mapsto u_\bullet(\widehat x_k),\quad
 \omega_E(\widehat x_\ell)\mapsto\omega_E(\widehat x_k).
\end{equation}
De \eqref{eq:pdfv2-cor-cell-frontera-aristas}--
\eqref{eq:pdfv2-cor-cell-frontera-celda} se obtiene
\begin{equation}
\label{eq:pdfv2-cor-collar-naturalidad}
\begin{aligned}
 \operatorname{Cell}(\tau)\partial
 &=\partial\operatorname{Cell}(\tau),\\
 \tau_VP_L&=P_{\tau_{\mathscr I}L}\tau_V,\qquad
 \tau_QB=B\tau_a,\qquad W_c\circ\tau_Q=W_c,\\
 \operatorname{Cell}(\tau)\vartheta_L
 &=\vartheta_{\tau_{\mathscr I}L}\operatorname{Cell}(\tau),\\
 (\tau_V\otimes\operatorname{Cell}_0(\tau))
 \eta^{\rm cell}_{\widehat x_\ell}
 &=\eta^{\rm cell}_{\widehat x_k}\tau_V.
\end{aligned}
\end{equation}
Identidad y composición se verifican en cada generador. Así
\(\operatorname{Cell}_{\rm TPK}\) es la completación celular funtorial de
la misma fibra APP--TRIT--TPK, no un portador añadido desde una publicación
posterior.

\subsection{Medida de incidencia, curvatura de acarreo y naturalidad adjunta}
En el desarrollo siguiente se identifican, exclusivamente como abreviaturas
de la completación común,
\[
 I_{\rm gau}:=\mathscr I_{\widehat x},\qquad
 \mathscr B_{\rm gau}:=\mathscr B_{\widehat x},\qquad
 Q_{\rm gau}:=Q_{\rm inc}(\widehat x).
\]
Con el promedio uniforme de conteo sobre \(Q_{\rm gau}\), la función
\(W_c\) de \eqref{eq:pdfv2-cor-wc} satisface
\begin{equation}
\label{eq:pdfv2-cor-wc-momentos}
 \mathbb E_QW_c=0,\qquad
 \mathbb E_QW_c^2=\frac29,\qquad
 \mathbb E_QW_c^3=\frac{2}{27}.
\end{equation}
Estas identidades son propiedades de la misma celda y del mismo acarreo, no
datos de una teoría exterior.
\begingroup
  \let\section\subsection
  \input{incorporaciones/v2_global_20260823/05c_incidencia_curvatura_correlativa}
\endgroup

La misma regla define su acción sobre una emisión elemental
\(\alpha:x\to y\). Como \(U_\alpha\) transporta la pareja \emph{ordenada}
de hojas, su mapa de etiquetas es
\(\overline U_\alpha e_{\Sigma,x}=e_{\Sigma,y}\) y
\(\overline U_\alpha e_{\Pi,x}=e_{\Pi,y}\). Su proyectivización
\(\alpha_{\mathscr I}[v]:=[\overline U_\alpha v]\) actúa en
\(\mathbf P^1(V_x)\) y, por tanto,
\begin{equation}
\label{eq:pdfv2-cor-cell-accion-emision}
 \begin{gathered}
 \operatorname{Cell}_{\rm TPK}(\alpha)v_\star(x)=v_\star(y),\qquad
 \operatorname{Cell}_{\rm TPK}(\alpha)v_L(x)
   =v_{\alpha_{\mathscr I}L}(y),\\
 \operatorname{Cell}_{\rm TPK}(\alpha)v_{\{L,L'\}}(x)
   =v_{\{\alpha_{\mathscr I}L,\alpha_{\mathscr I}L'\}}(y),\qquad
 \operatorname{Cell}_{\rm TPK}(\alpha)v_{L,\epsilon}(x)
   =v_{\alpha_{\mathscr I}L,\epsilon}(y),\\
 \operatorname{Cell}_{\rm TPK}(\alpha)u_L(x)
   =u_{\alpha_{\mathscr I}L}(y),\qquad
 \operatorname{Cell}_{\rm TPK}(\alpha)u_{L,\epsilon}(x)
   =u_{\alpha_{\mathscr I}L,\epsilon}(y),\\
 \operatorname{Cell}_{\rm TPK}(\alpha)u_{L'\mid L}(x)
   =u_{\alpha_{\mathscr I}L'\mid\alpha_{\mathscr I}L}(y),\qquad
 \operatorname{Cell}_{\rm TPK}(\alpha)\omega_{\{L,L'\}}(x)
   =\omega_{\{\alpha_{\mathscr I}L,
                 \alpha_{\mathscr I}L'\}}(y).
 \end{gathered}
\end{equation}
En la fibra de hojas y en su pantalla se tiene además
\begin{equation}
\label{eq:pdfv2-cor-cell-entrelaza-proyectores-pantalla}
 \overline U_\alpha P_L=P_{\alpha_{\mathscr I}L}\overline U_\alpha,
 \qquad
 \operatorname{Cell}_{\rm TPK}(\alpha)\vartheta_L
 =\vartheta_{\alpha_{\mathscr I}L}
  \operatorname{Cell}_{\rm TPK}(\alpha).
\end{equation}
Las fórmulas de frontera dan
\begin{equation}
\label{eq:pdfv2-cor-cell-accion-emision-natural}
 \partial\operatorname{Cell}_{\rm TPK}(\alpha)
 =\operatorname{Cell}_{\rm TPK}(\alpha)\partial,\qquad
 \operatorname{Cell}_{\rm TPK}(\beta\alpha)
 =\operatorname{Cell}_{\rm TPK}(\beta)
  \operatorname{Cell}_{\rm TPK}(\alpha).
\end{equation}
Así el mapa de complejos usado en el registro terminal es la acción libre
inducida por la misma emisión, no un símbolo añadido.

\subsection{Módulo de rutas, totalización y cono del retorno nonádico}
Sea \(J_d(x)\) la categoría completa de prefijos de profundidad \(d\) del
árbol común, \(\operatorname{Term}_d(x)\) su conjunto terminal y
\(\Gamma_{d,N}(\nu):=C_{d,N}(\nu)\) el complejo celular de
\eqref{eq:pdfv2-cor-cono-celda}. El módulo, la totalización y el cono que
siguen conservan simultáneamente todas las continuaciones; no seleccionan
una historia ni separan una construcción cohomológica del continuo común.
\begingroup
  \let\section\subsection
  \input{incorporaciones/v2_global_20260823/05d_totalizacion_cohomologica_correlativa}
\endgroup

El complejo local \eqref{eq:pdfv2-cor-complejo-local} y el vértice
\(y_\alpha\), extremo de esa misma extensión en el nervio, forman el pullback
aumentado. En él se definen
\begin{equation}
\label{eq:pdfv2-cor-zpm-h}
 \mathbf1_\alpha=(r,[y_\alpha]),\qquad
 z_-(\alpha)=(r,[y_\alpha]),
\end{equation}
\begin{equation}
\label{eq:pdfv2-cor-zplus-h}
 z_+(\alpha)=
 (r+3c_\alpha v_\alpha e+v_\alpha b,[y_\alpha]),
 \qquad
 h_*(\alpha)=(-v_\alpha f,0).
\end{equation}
La acción en generadores da, sin una existencia abstracta adicional,
\begin{equation}
\label{eq:pdfv2-cor-filler}
 \partial z_-=\partial z_+=0,\qquad
 \partial h_*=z_--z_+.
\end{equation}
El relleno correlativo, la celda y el collar conservan el mismo
\((\alpha,\operatorname{or}(\alpha),\kappa(\alpha),\partial\alpha,
\gamma\alpha,m(\gamma\alpha;m_0))\); sólo aplican lectores internos distintos
al registro común.

Para el mismo complejo finito basado se define, sin pedir balance ni
aciclicidad, la línea determinante total
\begin{equation}
\label{eq:pdfv2-cor-linea-determinante-total}
 \operatorname{Det}(C_{k,N}):=
 \bigotimes_q
 \left(\bigwedge^{\rm top}C_{k,N,q}\right)^{(-1)^q}.
\end{equation}
El refinamiento induce el morfismo determinantal del mapa de complejos y
transporta a la vez los módulos de homología, las formas normales de Smith y
las páginas de Bockstein como campos del registro. Ni la existencia de
\eqref{eq:pdfv2-cor-linea-determinante-total} ni su transporte exige que la
homología se anule. Si el cono del refinamiento es contractible y orientado,
el morfismo se trivializa canónicamente; en el caso general se conserva el
cono y no se lo sustituye por un isomorfismo ficticio.

\subsection{Reducción terminal, forma de Smith y torre de Bockstein}
Para cada prefijo \(y\), el acarreo no orientado y el coeficiente orientado
de frontera del registro común admiten los levantamientos
\begin{equation}
\label{eq:detint-c-v-lifts}
 \widetilde c_y,\widetilde v_y\in\mathbf Z,\qquad
 c_{y,n}:=\widetilde c_y\bmod3^n,\qquad
 v_{y,n}:=\widetilde v_y\bmod3^n.
\end{equation}
La salida determinantal no es una entrada ni una rama separada: es el sector
producido por la misma reducción terminal cuando se verifican las condiciones
tipadas
\begin{equation}
\label{eq:detint-dominio}
\begin{split}
 \mathcal C_{\rm cont}^{\rm det}:=\bigl\{x\in\mathcal C_{\rm cont}^{\rm disc}:{}&
 x\text{ sobrevive a toda profundidad};\\
 &\widehat P_x^{\rm red}\text{ es finito y basado};\\
 &\chi(\widehat P_x^{\rm red})=0;\\
 &H_*(\widehat P_x^{\rm red}\otimes_{\mathbf Z_3}K_3)=0;\\
 &\text{la reducción produce un bloque cuadrado estable bajo refinamiento}
 \bigr\}.
\end{split}
\end{equation}
\begingroup
  \let\section\subsection
  \input{incorporaciones/v2_global_20260823/05e_smith_bockstein_correlativo}
\endgroup

\section{Naturalidad del ensamblaje y compatibilidad bajo refinamiento}

Antes de formar el registro de horizonte, hacemos explícita la acción total
de las cinco operaciones aquí desarrolladas sobre \emph{la misma} arista
\(\mathfrak e_\alpha\). Si \(P_x,Q_x\) denotan las restricciones de
\eqref{eq:pdfv2-cor-sigma} a la fibra de \(x=s(\alpha)\), y análogamente en
\(y=t(\alpha)\), ponemos primero
\begin{equation}
\label{eq:pdfv2-cor-a-eta-emision}
 A_\alpha:=P_yU_\alpha P_x+Q_yU_\alpha Q_x,
 \qquad
 \eta_\alpha:=Q_yU_\alpha P_x+P_yU_\alpha Q_x.
\end{equation}
Las cinco lecturas son familias dependientes indexadas por la propia
\(\mathfrak e_\alpha\). El índice fija fuente, destino, ruta y ledger, pero
no se vuelve a copiar dentro del valor. Sus tipos son los registros
dependientes
\(\operatorname{LedgerEntry}:=\operatorname{cod}(\lambda)\) y
\begin{align*}
 \mathsf Y_{\rm sp}[\alpha]&:=
 \operatorname{Rec}[U:\mathsf H_x\to\mathsf H_y;
 \sigma_x\in\operatorname{End}\mathsf H_x;
 \sigma_y\in\operatorname{End}\mathsf H_y;A,\eta:\mathsf H_x\to\mathsf H_y],\\
 \mathsf Y_{\rm sol}[\alpha]&:=
 \operatorname{Rec}[b,b^+,b^-\in\operatorname{LedgerEntry};
 \operatorname{upd}:\mathsf{Mem}_x\to\mathsf{Mem}_y],\\
 \mathsf Y_{\rm gau}[\alpha]&:=
 \mathcal P_{\rm fin}\operatorname{Rec}[
 \operatorname{Cell}(\alpha);q\in Q_{\rm inc}(y);B;W_c(q)],\\
 \mathsf Y_{\rm coh}[\alpha]&:=
 \operatorname{Rec}[I_y,J_y,\kappa_{y,N},\delta_y,C_{y,N};
 F_\alpha:C_{y,N}\to C_{x,N}],\\
 \mathsf Y_{\rm det}[\alpha]&:=
 \operatorname{Rec}[\Psi_\alpha,K_\alpha,\mathbf1_\alpha,z_\pm,h_*;
 \operatorname{Det},H_*,\operatorname{Smith},\operatorname{Boc}].
\end{align*}
Por tanto el juicio común es
\[
 \mathfrak e_\alpha:\operatorname{Ext}_k\ \vdash\
 \mathcal O_\bullet(\mathfrak e_\alpha)
 \in\mathsf Y_\bullet[\alpha],
 \qquad
 \bullet\in\{\mathrm{sp,sol,gau,coh,det}\}.
\]
\begin{align}
 \mathcal O_{\rm sp}(\mathfrak e_\alpha)
 &:=
 \bigl(U_\alpha,\sigma_x,\sigma_y,
 A_\alpha,\eta_\alpha\bigr),
 \label{eq:pdfv2-cor-accion-sp}\\
 \mathcal O_{\rm sol}(\mathfrak e_\alpha)
 &:=
 \bigl(b_\alpha,b_\alpha^+,b_\alpha^-;\,
 (M^+,M^-,D,H)\mapsto
 (M^++b_\alpha^+,M^-+b_\alpha^-,
  D+b_\alpha,H+|b_\alpha|)\bigr),
 \label{eq:pdfv2-cor-accion-sol}\\
 \mathcal O_{\rm gau}(\mathfrak e_\alpha)
 &:=
 \left\{
 \bigl(\operatorname{Cell}_{\rm TPK}(\alpha),
 q,B,W_c(q)\bigr):
 q\in Q_{\rm inc}(y)
 \right\},
 \label{eq:pdfv2-cor-accion-gau}\\
 \mathcal O_{\rm coh}(\mathfrak e_\alpha)
 &:=
 \bigl(I_y,J_y,\kappa_{y,N},\delta_y,C_{y,N},
 F_\alpha:C_{y,N}\longrightarrow C_{x,N}\bigr),
 \label{eq:pdfv2-cor-accion-coh}\\
 \mathcal O_{\rm det}(\mathfrak e_\alpha)
 &:=
 \bigl(\Psi_\alpha,K_\alpha,\mathbf1_\alpha,z_\pm(\alpha),
 h_*(\alpha),\operatorname{Det}(C_{y,N}),H_*(C_{y,N}),
 \operatorname{Smith}(C_{y,N}),\operatorname{Boc}(C_{y,N})\bigr).
 \label{eq:pdfv2-cor-accion-det}
\end{align}

La lectura arquimediana interna es la primera componente de esta acción
simultánea; no introduce todavía ninguna forma clásica:
\begin{equation}
\label{pII-full:eq:continuo-lector-arquimediano}
 \mathcal L_{\rm arq}^{\rm int}(\mathfrak e_\alpha)
 :=\mathcal O_{\rm sp}(\mathfrak e_\alpha).
\end{equation}
Para cualquier valor de esas familias, incluido cada elemento de la
multisección \(\mathcal O_{\rm gau}\), su soporte es el índice dependiente
\begin{equation}
\label{eq:pdfv2-cor-soporte-dependiente}
 \begin{aligned}
 \mathsf B_\alpha&:=
 \{(s(\alpha),t(\alpha),\gamma\alpha,
          \operatorname{Led}(\gamma\alpha))\},\\
 \operatorname{supp}_\alpha:\mathsf Y_\bullet[\alpha]&\longrightarrow
 \mathsf B_\alpha,\qquad
 o\longmapsto(s(\alpha),t(\alpha),\gamma\alpha,
          \operatorname{Led}(\gamma\alpha)).
 \end{aligned}
\end{equation}
Por tanto quedan tipados, sin una proyección que recupere la arista,
\begin{equation}
\label{eq:pdfv2-cor-soporte-cinco-lecturas}
 \begin{gathered}
 s_\alpha(\mathcal O_\bullet):=s(\alpha),\qquad
 t_\alpha(\mathcal O_\bullet):=t(\alpha),\\
 \gamma_\alpha(\mathcal O_\bullet):=\gamma\alpha,
 \qquad
 \operatorname{Led}_\alpha(\mathcal O_\bullet)
 :=\operatorname{Led}(\gamma\alpha).
 \end{gathered}
\end{equation}
Aquí \(\operatorname{Cell}_{\rm TPK}(\alpha)\) es el mapa de complejos
\eqref{eq:pdfv2-cor-cell-accion-emision} inducido por el transporte de la
misma fibra de hojas, y \(F_\alpha\) es
\eqref{eq:pdfv2-cor-refinamiento-celda} para esa extensión. En particular,
\(\mathcal O_{\rm gau}\) conserva la multisección completa de \(q\)'s; no
elige el origen ni una órbita.

Las cinco fórmulas anteriores constituyen una sola acción dependiente
\begin{equation}
\label{eq:pdfv2-cor-accion-unica-dependiente}
 \mathcal O(\mathfrak e_\alpha)
 :=\left\langle
 \mathcal O_{\rm sp},\mathcal O_{\rm sol},\mathcal O_{\rm gau},
 \mathcal O_{\rm coh},\mathcal O_{\rm det};
 \operatorname{Cpl}_\alpha
 \right\rangle,
\end{equation}
donde \(\operatorname{Cpl}_\alpha\) es el certificado formado por las
identidades
\begin{equation}
\label{eq:pdfv2-cor-acoplamientos-unicos}
 \begin{gathered}
 \operatorname{supp}_\alpha(\mathcal O_{\rm sp})
 =\operatorname{supp}_\alpha(\mathcal O_{\rm sol})
 =\operatorname{supp}_\alpha(\mathcal O_{\rm gau})
 =\operatorname{supp}_\alpha(\mathcal O_{\rm coh})
 =\operatorname{supp}_\alpha(\mathcal O_{\rm det}),\\
 \mathsf Q(t(\alpha))=\sigma_9\mathsf Q(s(\alpha)),\quad
 g(t(\alpha))=g(s(\alpha))+[1]_9,\quad
 q^{(9)}(t(\alpha))=q^{(9)}(s(\alpha))+\chi_9(g(s(\alpha))),\\
 A_\alpha^*A_\alpha+\eta_\alpha^*\eta_\alpha=U_\alpha^*U_\alpha,\qquad
 (M^+,M^-,D,H)'=(M^++b_\alpha^+,M^-+b_\alpha^-,
 D+b_\alpha,H+|b_\alpha|),\\
 Q_{\rm inc}=C^2(\operatorname{Cell}_{\rm TPK};\mathbf F_3),\qquad
 sB=0,\qquad W_c(q+Ba)=W_c(q),\\
 C_{y,N}\text{ es el mismo complejo en }
 \mathcal O_{\rm coh}\text{ y }\mathcal O_{\rm det},\qquad
 \partial h_*=z_--z_+,\qquad
 K_{\beta\alpha}=K_\beta+\Psi_\beta K_\alpha.
 \end{gathered}
\end{equation}
El símbolo \(\bullet\) recorre las cinco lecturas de
\eqref{eq:pdfv2-cor-accion-unica-dependiente}; todas conservan la misma
fuente, destino, ruta y ledger.

Reunimos, sin separar el estado, las operaciones construidas hasta aquí:
\begin{equation}
\label{eq:pdfv2-cor-constructor-unico}
 \begin{aligned}
 \mathfrak D_{k,N}(\widehat x)
 &:=\bigl(\{\mathfrak e_\alpha\}_\alpha;\\
 &\quad U_k,\sigma_k,P_k,Q_k,A_k,\eta_k;\\
 &\quad b_\alpha,b_\alpha^\pm,M_k^\pm,D_k^0,H_k^0,
 \kappa_k^{\rm can};\\
 &\quad V_{\rm TPK},\mathscr I,\mathscr E,\mathscr B,B;\\
 &\quad \operatorname{Cell}_{\rm TPK}(\widehat x),Q_{\rm inc}(\widehat x),W_c;\\
 &\quad I_k,J_k,\kappa_{k,N},\delta,C_{k,N};\\
 &\quad c_\alpha,v_\alpha,\Psi_\alpha,K_\alpha,
 \mathbf1_\alpha,z_\pm(\alpha),h_*(\alpha);\\
 &\quad
 \operatorname{Det}(C_{k,N}),H_*(C_{k,N}),
 \operatorname{Smith}(C_{k,N}),\operatorname{Boc}(C_{k,N});\\
 &\quad \{\mathcal O(\mathfrak e_\alpha)\}_\alpha
 \bigr).
 \end{aligned}
\end{equation}
El punto y coma sólo separa niveles de lectura de una misma tupla dependiente;
no define estructuras independientes. El estado de entrada no se guarda como
primera coordenada ni se recupera mediante una proyección. Cada operación se
demuestra por su acción en los generadores \(\mathfrak e_\alpha\), que
conservan el mismo soporte, la misma ruta y el mismo ledger.

Para \(\ell\geq k\) y \(N'\geq N\), el transporte se define componente a
componente por truncación de historias, reducción de coeficientes y
precomposición de puertos. La afirmación de naturalidad que debe probarse es
una sola:
\begin{equation}
\label{eq:pdfv2-cor-naturalidad-unica}
 u_{(\ell,N'),(k,N)}\mathfrak D_{\ell,N'}(\widehat x_\ell)
 =\mathfrak D_{k,N}(\tau_{\ell,k}\widehat x_\ell).
\end{equation}
La naturalidad se formula sobre la historia generada, no fingiendo que una
arista fina aislada deba seguir siendo arista después de truncar sus dos
extremos. Para
\(\gamma=(\varepsilon_1,\ldots,\varepsilon_\ell)\), definimos
\[
 \mathsf E_{\leq\ell}(\gamma)
 :=(\mathfrak e_{\varepsilon_i})_{1\leq i\leq\ell},
 \qquad
 \mathcal O_{\leq\ell}(\gamma)
 :=(\mathcal O(\mathfrak e_{\varepsilon_i}))_{1\leq i\leq\ell}.
\]
Los dos transportes son los mapas de prefijo explícitos
\begin{align}
 u^{\mathfrak e}_{\ell,k}:
 \mathsf E_{\leq\ell}(\gamma)&\longrightarrow
 \mathsf E_{\leq k}(\gamma),&
 (\mathfrak e_{\varepsilon_i})_{i\leq\ell}
 &\longmapsto(\mathfrak e_{\varepsilon_i})_{i\leq k},
 \label{eq:pdfv2-cor-transporte-arista}\\
 u^{\mathcal O}_{\ell,k}:
 \mathcal O_{\leq\ell}(\gamma)&\longrightarrow
 \mathcal O_{\leq k}(\gamma),&
 (\mathcal O(\mathfrak e_{\varepsilon_i}))_{i\leq\ell}
 &\longmapsto
 (\mathcal O(\mathfrak e_{\varepsilon_i}))_{i\leq k}.
 \label{eq:pdfv2-cor-transporte-accion}
\end{align}
En cada entrada conservada, la componente \(\mathrm{sp}\) entrelaza
\(U,\sigma,P,Q,A,\eta\); la componente \(\mathrm{sol}\) aplica la esperanza
condicional a \(b,b^\pm,\kappa\); la componente \(\mathrm{gau}\) aplica
\(\operatorname{Cell}(\tau)\oplus\tau_Q\); la componente \(\mathrm{coh}\)
precompone los puertos con \(F\); y la componente \(\mathrm{det}\) aplica
\(\Psi,K,H_*,\operatorname{Smith},\operatorname{Boc}\) al mismo mapa de
complejos. El prefijo actúa simultáneamente sobre esas cinco componentes.
Sus componentes ya descargadas son
\begin{equation}
\label{eq:pdfv2-cor-naturalidad-componentes}
 \begin{gathered}
 U_{m,k}=U_{m,\ell}U_{\ell,k}\quad(m\geq\ell\geq k),\qquad
 E_kb_f^\pm=b_c^\pm+\kappa_k^{\rm can},\\
 r_{N',N}\kappa_{k,N'}=
 \kappa_{k,N}(r_{N',N}\oplus r_{N',N}),\\
 \Psi_\beta\Psi_\alpha=\Psi_{\beta\alpha},
 \qquad K_{\beta\alpha}=K_\beta+\Psi_\beta K_\alpha,\\
 u^{\mathcal O}_{\ell,k}\mathcal O_{\leq\ell}(\gamma)
 =\mathcal O_{\leq k}(\gamma)
 =\operatorname{map}(\mathcal O)\bigl(u^{\mathfrak e}_{\ell,k}
                    \mathsf E_{\leq\ell}(\gamma)\bigr).
 \end{gathered}
\end{equation}
La primera igualdad usa la multiplicación de los pesos condicionales sobre el mismo
sistema de cilindros comunes; las restantes son las fórmulas explícitas
anteriores. Ninguna proyección visible
puede usarse para reconstruir o seleccionar \(\widehat x\).

Para completar el último componente, denotemos por
\(F_{(\ell,N'),(k,N)}:C_{\ell,N'}\to C_{k,N}\) el mapa de complejos de
\eqref{eq:pdfv2-cor-refinamiento-celda}. Su cono se conserva como parte de la
transición. La functorialidad de homología y de la sucesión de Bockstein da
\begin{equation}
\label{eq:pdfv2-cor-naturalidad-homologia-bockstein}
 H_*(F_{m,k})=H_*(F_{\ell,k})H_*(F_{m,\ell}),
 \qquad
 \operatorname{Boc}(F_{m,k})
 =\operatorname{Boc}(F_{\ell,k})
  \operatorname{Boc}(F_{m,\ell}).
\end{equation}
La línea determinante se transporta mediante la identidad del triángulo
exacto
\begin{equation}
\label{eq:pdfv2-cor-naturalidad-determinante}
 \operatorname{Det}(C_{k,N})
 \simeq
 \operatorname{Det}(C_{\ell,N'})\otimes
 \operatorname{Det}\!\left(\operatorname{Cone}
 F_{(\ell,N'),(k,N)}\right).
\end{equation}
Esta fórmula registra también los refinamientos que añaden homología. La
lectura de Smith se conserva sin fingir invariancia: a cada transición se le
asigna el certificado total
\begin{equation}
\label{eq:pdfv2-cor-transicion-smith}
 \mathcal S(F):=
 \bigl(\operatorname{Smith}(C_{\ell,N'}),
       \operatorname{Smith}(C_{k,N}),
       \operatorname{Smith}(\operatorname{Cone}F)\bigr),
\end{equation}
que determina literalmente qué sumandos se conservan, nacen o mueren.

La acción de una simetría coherente \(g\) en el árbol común induce
isometrías de permutación \(\widetilde g_k\) en
\(\mathscr H_k\). La biyección de hijos y la conservación de los pesos de
\eqref{eq:continuo-comun-naturalidad-medida} prueban en los vectores básicos
\begin{equation}
\label{eq:pdfv2-cor-naturalidad-transporte-hojas}
 \widetilde g_{k+1}U_k=U_k\widetilde g_k,
 \qquad
 \widetilde g_k\sigma_k=\sigma_k\widetilde g_k,
 \qquad
 \widetilde g_kP_k=P_k\widetilde g_k,
 \qquad
 \widetilde g_kQ_k=Q_k\widetilde g_k.
\end{equation}

\begin{theorem}[Totalidad y naturalidad del constructor correlativo]
\label{thm:pdfv2-cor-totalidad-naturalidad}
Para todo estado total de \eqref{eq:continuo-comun-estado}, todo horizonte
finito y todo nivel de coeficientes, la expresión
\eqref{eq:pdfv2-cor-constructor-unico} está definida. Sus cinco
características son operaciones simultáneas sobre el mismo registro de
extensión y satisfacen la única identidad
\eqref{eq:pdfv2-cor-naturalidad-unica}. El collar, el complejo y la línea
determinante permanecen no vacíos aun cuando sus lectores posteriores no
sean acíclicos ni invertibles.
\end{theorem}

\begin{proof}
La no vacuidad del árbol y de sus puertos es
\cref{prop:continuo-comun-arbol-finito-no-vacio}; la medida y la expectativa
son totales por \cref{thm:continuo-comun-medida-canonica}. Las fórmulas
  \eqref{eq:pdfv2-cor-u-comun}--\eqref{eq:pdfv2-cor-memorias} definen transporte,
polaridad, balance y memoria para todos los generadores. Las cuatro
  direcciones de \eqref{eq:pdfv2-cor-cuatro-rectas-tpk} son distintas por
  \eqref{eq:pdfv2-cor-proyectores-no-colapso}; la completación celular
  \cref{def:pdfv2-cor-completacion-celular-tpk} contiene sus dos órdenes de
  camino y seis celdas, mientras
  \eqref{eq:pdfv2-cor-collar-naturalidad} prueba su naturalidad. La
sucesión celular es exacta por \eqref{eq:pdfv2-cor-sucesion-celda}, el
complejo local cumple \(\partial^2=0\) y el relleno correlativo es explícito por
\eqref{eq:pdfv2-cor-filler}. Toda construcción finita basada posee la línea
\eqref{eq:pdfv2-cor-linea-determinante-total}; homología, Bockstein y el
certificado de Smith se transportan por
\eqref{eq:pdfv2-cor-naturalidad-homologia-bockstein}--
\eqref{eq:pdfv2-cor-transicion-smith}. Finalmente, las identidades se
verifican en cada generador \(\mathfrak e_\alpha\) mediante
\eqref{eq:pdfv2-cor-naturalidad-componentes},
\eqref{eq:pdfv2-cor-collar-naturalidad} y
\eqref{eq:pdfv2-cor-naturalidad-transporte-hojas}; por linealidad y por la
composición functorial de rutas se obtiene
\eqref{eq:pdfv2-cor-naturalidad-unica} para el registro entero.
\end{proof}

\section{Estratificación de las afirmaciones del constructor correlativo: identidades internas, deducciones relativas e interfaces físicas}

\paragraph{No vacuidad e identidades exactas del constructor correlativo}
La no vacuidad importada de \eqref{eq:continuo-comun-out-finito} y las
identidades
\eqref{eq:pdfv2-cor-b-identidades},
\eqref{eq:pdfv2-cor-sucesion-celda},
\eqref{eq:pdfv2-cor-sb-cero},
\eqref{eq:pdfv2-cor-u-gram},
\eqref{eq:pdfv2-cor-balance-hojas},
\eqref{eq:pdfv2-cor-jensen-hojas},
\eqref{eq:pdfv2-cor-dos-cero},
\eqref{eq:pdfv2-cor-cociclo-matricial} y
\eqref{eq:pdfv2-cor-filler} se deducen de las definiciones y de las relaciones
APP--TRIT--TPK exhibidas. Su estatuto es
\texttt{EXACTO\_INTERNO}.

\paragraph{Hipótesis de partida para la isometría de \(U_k\), la sucesión celular y la composición \(\Psi/K\)}
La construcción de \(U_k\) usa como estructura de partida la finitud de las
fibras, la supervivencia a todo horizonte, el truncamiento funcional y los
transportes de las dos hojas ya producidos por TPK\@. La igualdad
\(U_k^*U_k=I\) exige además verificar \(U_\alpha^*U_\alpha=I\) para cada
emisión; sin esa verificación permanece la identidad de Gram
\eqref{eq:pdfv2-cor-u-gram}, que sí es incondicional. La sucesión celular usa
los dos puertos APP ordenados. La composición \(\Psi/K\) usa carry entero,
orientación y frontera. Bajo esas premisas enumeradas, las conclusiones
indicadas son afirmativas.

\paragraph{Datos adicionales para realizaciones isométricas, retractos balanceados y trivialización de la línea determinante}
El constructor interno queda completo con el Gram positivo de
\eqref{eq:pdfv2-cor-u-gram}, el espacio total \(Q_{\rm inc}\), la sucesión
celular y la línea determinante con su cono de transición. Una publicación
especializada puede exigir datos adicionales: la isometría de cada
\(U_\alpha\), una identificación numérica concreta del defecto APP, un
retracto balanceado y acíclico o la trivialización de una línea
determinante. Esos datos son entrelazadores del codominio publicado; no
forman parte del dominio ni de la existencia de
\(\mathfrak D_{k,N}\). Su ausencia impide exclusivamente la promoción que
los mencione.

\begin{criterion}[Falsadores de la correlación intrínseca]
\label{crit:pdfv2-cor-falsadores}
La construcción conjunta falla en el primer nivel en que ocurra cualquiera de
los hechos siguientes:
\begin{enumerate}
 \item una operación usa una extensión, orientación, carry, memoria, frontera
 o ruta diferente de la registrada en \(\mathfrak e_\alpha\);
 \item se escoge un hijo de
 \(\operatorname{Ch}_k(\widehat x_k)\) y se eliminan los restantes;
 \item el peso de \eqref{eq:pdfv2-cor-peso-conteo} no suma uno, falla
 \eqref{eq:pdfv2-cor-u-gram}, o se declara \(U_k^*U_k=I\) sin probar
 \(U_\alpha^*U_\alpha=I\) para toda emisión;
 \item se identifica \(\kappa_k^{\rm can}\) con el carry entero
 \(\kappa_\alpha\);
 \item falla \(sB=0\), se contrae
 \(Q_{\rm inc}(\widehat x)\) a su origen, se altera la acción
 \(q\mapsto q+Ba\), o se identifican los dos órdenes de camino antes de
 aplicar el lector;
 \item un refinamiento rompe
 \eqref{eq:pdfv2-cor-refinamiento-cadena} o el cociclo
 \eqref{eq:pdfv2-cor-cociclo-matricial};
 \item falla \(\partial^2=0\), \(z_\pm\) no son ciclos o
 \(\partial h_*\neq z_--z_+\);
 \item se afirma un escalar determinantal no nulo sin el retracto, el
 balance y la aciclicidad exigidos por esa publicación;
 \item se presenta cualquiera de las cinco operaciones como un objeto
 separado del constructor \(\mathfrak D_{k,N}\).
\end{enumerate}
En todos esos casos el diagnóstico es: \emph{objeto común sustituido;
conclusión no transferible}.
\end{criterion}
